% Source: Wald GR, physical PDF pp. 436--440
%         (printed pp. 423--427).
% Coverage: Appendix A, Topological Spaces; Theorems A.1--A.9.
% Figures/tables/footnotes: none.
% Status: source-reviewed and compile-clean.
% Uncertainties: none.

\section{Topological Spaces}

In mathematics, many different proofs and arguments are based on the same or a
similar set of ideas. An important goal of mathematics is to isolate the key ideas in
as general a form as possible, since if one can derive general results about these ideas
in the abstract, one then may apply them to the cases of interest without duplication
of effort. The notion of a topological space provides a beautiful illustration of this
program. One starts with a very simple, abstract definition and ends up proving many
powerful results that have a very wide range of application. Our main interest here
in topological spaces arises from the fact that in general relativity spacetime has the
structure of a topological space. Indeed, spacetime has a great deal more structure,
and we circumvented the mention of topological spaces in Chapter~2 by defining the
notion of a manifold directly. However, the purely topological arguments and results
discussed in this section play an important role in many of the constructions and
proofs of Chapters~8 and~9. Therefore, we collect in this appendix many of the key
definitions and theorems concerning topological spaces.

A \emph{topological space} \((X,\mathcal T)\) consists of a set \(X\) together with a
collection \(\mathcal T\) of subsets of \(X\) satisfying the following three properties:

\begin{enumerate}[label=\textup{(\arabic*)},leftmargin=2.2em]
\item
The union of an arbitrary collection of subsets, each of which is in \(\mathcal T\), is
in \(\mathcal T\): If \(O_\alpha\in\mathcal T\) for all \(\alpha\), then
\(\bigcup_\alpha O_\alpha\in\mathcal T\).

\item
The intersection of a finite number of subsets in \(\mathcal T\) is in \(\mathcal T\):
If \(O_1,\ldots,O_n\in\mathcal T\), then
\[
  \bigcap_{i=1}^{n}O_i\in\mathcal T.
\]

\item
The entire set \(X\) and the empty set \(\emptyset\) are in \(\mathcal T\).
\end{enumerate}

\(\mathcal T\) is referred to as a \emph{topology} on \(X\), and subsets of \(X\)
which are listed in the collection \(\mathcal T\) are called \emph{open sets}.

Any set \(X\) can easily be made into a topological space by taking
\(\mathcal T=\{\text{all subsets of \(X\)}\}\) (the \emph{discrete topology}) or by
taking \(\mathcal T=\{X,\emptyset\}\) (the \emph{indiscrete topology}). A much more
interesting example of a topological space is obtained by taking \(X=\mathbb R\), the
set of real numbers, and defining \(\mathcal T\) to consist of all subsets of
\(\mathbb R\) which can be expressed as unions of open intervals \((a,b)\). (Thus,
with this choice of \(\mathcal T\) on \(\mathbb R\) an open interval is an open set;
historically, this example is the reason why the terminology \lq\lq open set\rq\rq\ is
used in the discussion of an abstract topological space.) More generally, for any
metric space, the collection of all subsets which can be expressed as unions of open
balls yields a topology.

If \((X,\mathcal T)\) is a topological space and \(A\) is any subset of \(X\), we may
make \(A\) itself into a topological space by defining the topology, \(\mathcal S\), on
\(A\) to consist of all subsets of \(A\) which can be expressed as intersections of
elements of \(\mathcal T\) with \(A\),
\[
  \mathcal S=\{U\mid U=A\cap O,\ O\in\mathcal T\}.
\]
\(\mathcal S\) is called the \emph{induced} (or \emph{relative}) topology.

If \((X_1,\mathcal T_1)\) and \((X_2,\mathcal T_2)\) are topological spaces, we can
make the product space
\[
  X_1\times X_2=\{(x_1,x_2)\mid x_1\in X_1,\ x_2\in X_2\}
\]
into a topological space \((X_1\times X_2,\mathcal T)\) by defining \(\mathcal T\) to
consist of all subsets of \(X_1\times X_2\) which can be expressed as unions of sets
of the form \(O_1\times O_2\) with \(O_1\in\mathcal T_1\) and
\(O_2\in\mathcal T_2\). \(\mathcal T\) is called the \emph{product topology}, and using
the above definition of a topology on \(\mathbb R\), we can use the construction of
the product topology to define a topology on \(\mathbb R^n\). The topology we get is
the same one as would be obtained by directly defining \(\mathcal T\) to consist of
all subsets of \(\mathbb R^n\) which can be expressed as unions of open balls.

If \((X,\mathcal T)\) and \((Y,\mathcal S)\) are topological spaces, a map
\(f:X\to Y\) is said to be \emph{continuous} if the inverse image,
\(f^{-1}(O)=\{x\in X\mid f(x)\in O\}\), of every open set \(O\) in \(Y\) is an open
set in \(X\). For functions from \(\mathbb R\) into \(\mathbb R\), it is easy to verify
that with the above definition of a topology on \(\mathbb R\), this definition of
continuity is equivalent to the usual \(\epsilon\)--\(\delta\) definition. If \(f\) is
continuous, one-to-one, onto, and its inverse is continuous, \(f\) is called a
\emph{homeomorphism} and \((X,\mathcal T)\) and \((Y,\mathcal S)\) are said to be
homeomorphic. Homeomorphic topological spaces have identical topological properties.

If \((X,\mathcal T)\) is a topological space, a subset \(C\) of \(X\) is said to be
\emph{closed} if its complement
\(X-C=\{x\in X\mid x\notin C\}\) is open. Thus, for example, a closed interval
\([a,b]\) of \(\mathbb R\) (with the standard topology on \(\mathbb R\)) is a closed
set. From the topological space axioms it follows immediately that the intersection
of an arbitrary collection of closed sets is closed and the union of a finite number of
closed sets is closed. Note that a subset may be neither open nor closed (e.g., the
half-open interval \([a,b)\) in \(\mathbb R\)) or may be both open and closed (as are
all subsets in the discrete topology). Indeed, the possibility of having subsets which
are both open and closed gives rise to a topological definition of connectedness: A
topological space \((X,\mathcal T)\) is said to be \emph{connected} if the only
subsets which are both open and closed are the entire space \(X\) and the empty set
\(\emptyset\). \(\mathbb R^n\) with the standard topology defined above is connected.

If \((X,\mathcal T)\) is a topological space and \(A\) is an arbitrary subset of \(X\),
the \emph{closure}, \(\overline A\), of \(A\) is defined as the intersection of all
closed sets containing \(A\). Clearly, \(\overline A\) is closed, contains \(A\), and
equals \(A\) if and only if \(A\) is closed. The \emph{interior} of \(A\) is defined as
the union of all open sets contained within \(A\). Clearly the interior of \(A\) is
open, is contained in \(A\), and equals \(A\) if and only if \(A\) is open. The
\emph{boundary} of \(A\), denoted \(\partial A\), consists of all points which lie in
\(\overline A\) but not in the interior of \(A\).

A topological space \((X,\mathcal T)\) is said to be \emph{Hausdorff} if for each pair
of distinct points \(p,q\in X\), \(p\ne q\), one can find open sets
\(O_p,O_q\in\mathcal T\) such that \(p\in O_p\), \(q\in O_q\), and
\(O_p\cap O_q=\emptyset\). It is easy to check that \(\mathbb R^n\) with the standard
topology is Hausdorff.

One of the most powerful notions in topology is that of compactness, which is
defined as follows. If \((X,\mathcal T)\) is a topological space and \(A\) is a subset
of \(X\), a collection \(\{O_\alpha\}\) of open sets is said to be an \emph{open cover}
of \(A\) if the union of these sets contains \(A\). A subcollection of the sets
\(\{O_\alpha\}\) which also covers \(A\) is referred to as a \emph{subcover}. The set
\(A\) is said to be \emph{compact} if every open cover of \(A\) has a finite subcover
(i.e., a subcover consisting of only a finite number of sets). Thus, for example, in
any topological space a set consisting of a single point is compact. On the other
hand, the open interval \((0,1)\) in \(\mathbb R\) (with the standard topology) is not
compact since the sets \(O_n=(1/n,1)\) for \(n=2,3,\ldots\), yield an open cover of
\((0,1)\) which admits no finite subcover.

The following theorems describe the implications of compactness and show the
utility of this notion. Proofs of the theorems can be found in nearly any textbook on
topology (e.g., Hocking and Young 1961; Kelley 1955).

Perhaps the most important theorem concerning compact subsets of \(\mathbb R\) is
the Heine--Borel theorem:

\begin{theorem}[Heine--Borel]\label{thm:A.1}
A closed interval \([a,b]\) of real numbers is compact (with the standard topology
on \(\mathbb R\)).
\end{theorem}

The general relation between compact and closed sets is described by the following
two theorems, the proofs of which are straightforward:

\begin{theorem}\label{thm:A.2}
Let \((X,\mathcal T)\) be Hausdorff and let \(A\subseteq X\) be compact. Then \(A\)
is closed.
\end{theorem}

\begin{theorem}\label{thm:A.3}
Let \((X,\mathcal T)\) be compact and let \(A\subseteq X\) be closed. Then \(A\) is
compact.
\end{theorem}

Combining the above three theorems, we arrive at the following strengthened
statement on the compactness of subsets of \(\mathbb R\):

\begin{theorem}\label{thm:A.4}
A subset \(A\) of real numbers is compact if and only if it is closed and bounded.
\end{theorem}

The property of compactness is easily proven to be preserved under continuous maps.
We have:

\begin{theorem}\label{thm:A.5}
Let \((X,\mathcal T)\) and \((Y,\mathcal S)\) be topological spaces. Suppose
\((X,\mathcal T)\) is compact and \(f:X\to Y\) is continuous. Then
\(f(X)=\{y\in Y\mid y=f(x)\}\) is compact.
\end{theorem}

Because of the properties of compact subsets of \(\mathbb R\) given by
Theorem~\ref{thm:A.4}, we have as a corollary

\begin{theorem}\label{thm:A.6}
A continuous function from a compact topological space into \(\mathbb R\) is bounded
and attains its maximum and minimum values.
\end{theorem}

The following theorem gives an immediate extension of results on compactness for
\(\mathbb R\) to results for \(\mathbb R^n\).

\begin{theorem}[Tychonoff theorem]\label{thm:A.7}
Let \((X_1,\mathcal T_1)\) and \((X_2,\mathcal T_2)\) be compact topological spaces.
Then the product space \(X_1\times X_2\) is compact in the product topology.
\end{theorem}

Theorem~\ref{thm:A.7} can be generalized to apply to the product of infinitely many
topological spaces, but the axiom of choice is needed for this generalization.

A corollary of this and the above theorems is

\begin{theorem}\label{thm:A.8}
A subset, \(A\), of \(\mathbb R^n\) is compact if and only if it is closed and bounded.
\end{theorem}

Thus, for example, the \(n\)-dimensional sphere \(S^n\) (defined as the set of points
in \(\mathbb R^{n+1}\) satisfying \(x_1^2+\cdots+x_{n+1}^2=1\)) in the induced
topology is compact, since it is easily seen to be a closed and bounded subset of
\(\mathbb R^{n+1}\).

A further notion we shall need is that of convergence of sequences. A sequence
\(\{x_n\}\) of points in a topological space \((X,\mathcal T)\) is said to converge to
point \(x\) if given any open neighborhood \(O\) of \(x\) (i.e., an open set \(O\)
containing \(x\)) there is an \(N\) such that \(x_n\in O\) for all \(n>N\). The point
\(x\) is said to be the limit of this sequence. It is easy to check that for
\(\mathbb R\) (with the standard topology) this agrees with the usual definition of
convergence. A point \(y\in X\) is said to be an \emph{accumulation point} (or
\emph{limit point}) of \(\{x_n\}\) if every open neighborhood of \(y\) contains
infinitely many points of the sequence. If \(\{x_n\}\) converges to \(x\), then \(x\)
clearly also is an accumulation point of the sequence. However, in a general
topological space, if \(y\) is an accumulation point of \(\{x_n\}\), it may not even
be possible to find a subsequence \(\{y_n\}\) of points of the sequence \(\{x_n\}\)
such that \(\{y_n\}\) converges to \(y\). However, the extraction of a subsequence
convergent to \(y\) will always be possible if \((X,\mathcal T)\) is \emph{first
countable}, that is, if for each \(p\in X\) there is a countable collection
\(\{O_n\}\) of open sets such that every open neighborhood, \(O\), of \(p\) contains
at least one member of this collection. \(\mathbb R^n\) with the standard topology is
first countable; indeed, it satisfies the stronger requirement of \emph{second
countability}: there is a countable collection of open sets such that every open set
can be expressed as a union of sets in this collection. For \(\mathbb R^n\), the open
balls with rational radii centered on points with rational coordinates compose such a
countable collection of open sets.

An important relation between compactness and convergence of sequences is expressed
by the Bolzano--Weierstrass theorem:

\begin{theorem}[Bolzano--Weierstrass theorem]\label{thm:A.9}
Let \((X,\mathcal T)\) be a topological space and let \(A\subseteq X\). If \(A\) is
compact, then every sequence \(\{x_n\}\) of points in \(A\) has an accumulation point
lying in \(A\). Conversely, if \((X,\mathcal T)\) is second countable and every
sequence in \(A\) has an accumulation point in \(A\), then \(A\) is compact. Thus, in
particular, if \((X,\mathcal T)\) is second countable, \(A\) is compact if and only if
every sequence in \(A\) has a convergent subsequence whose limit lies in \(A\).
\end{theorem}

Finally, we define the notion of paracompactness, a property which manifolds are
required to satisfy in order to prevent them from being \lq\lq too large.\rq\rq\
Let \((X,\mathcal T)\) be a topological space and let \(\{O_\alpha\}\) be an open
cover of \(X\). An open cover \(\{V_\beta\}\) is said to be a \emph{refinement} of
\(\{O_\alpha\}\) if for each \(V_\beta\) there exists an \(O_\alpha\) such that
\(V_\beta\subseteq O_\alpha\). The cover \(\{V_\beta\}\) is said to be \emph{locally
finite} if each \(x\in X\) has an open neighborhood \(W\) such that only finitely
many \(V_\beta\) satisfy \(W\cap V_\beta\ne\emptyset\). The topological space
\((X,\mathcal T)\) is said to be \emph{paracompact} if every open cover
\(\{O_\alpha\}\) of \(X\) has a locally finite refinement \(\{V_\beta\}\).

It is not difficult to show (see, e.g., Hocking and Young 1961) that any Hausdorff
topological space which is locally compact (i.e., such that every point has an open
neighborhood with compact closure) and which can be expressed as a countable union
of compact subsets is paracompact. Thus, \(\mathbb R^n\), \(S^m\) and their products
are easily verified to be paracompact. Indeed, it is not easy to construct examples of
topological spaces which satisfy all the requirements for a manifold but are not
paracompact; the \lq\lq long line\rq\rq\ (see Hocking and Young 1961) is perhaps the
simplest example, although the axiom of choice is required to define it.

For a manifold, \(M\), paracompactness has a number of important consequences. It
can be shown (see Kobayashi and Nomizu 1963) that paracompactness implies that
(1) \(M\) admits a Riemannian metric and (2) \(M\) is second countable. This latter
result implies, incidentally, that we can cover \(M\) by a locally finite, countable
family of charts \((\psi_i,O_i)\) with each \(\overline{O_i}\) compact. Conversely,
if \(M\) satisfies all the requirements of a manifold (see Chapter~2), then either of
properties (1) or (2) implies that \(M\) is paracompact.

Probably the most important consequence of paracompactness for a manifold is the
existence of a partition of unity. Given a locally finite open cover \(\{O_\alpha\}\)
of \(M\), a \emph{partition of unity} subordinate to \(\{O_\alpha\}\) is a collection
of smooth functions \(\{f_\alpha\}\) such that (i) the support of \(f_\alpha\) (i.e.,
the closure of the set where \(f_\alpha\) is nonvanishing) is contained within
\(O_\alpha\), (ii) \(0\leq f_\alpha\leq1\), and (iii) \(\sum_\alpha f_\alpha=1\).
[Since only finitely many \(f_\alpha\) are nonvanishing at any point, the sum in
property (iii) is only a finite sum.] It can be shown (Kobayashi and Nomizu 1963)
that every locally finite open cover \(\{O_\alpha\}\) of \(M\) such that each
\(\overline{O_\alpha}\) is compact admits a subordinate partition of unity. The
existence of a partition of unity allows us to globalize many local results. For
example, we can prove the above mentioned result that a paracompact manifold admits
a Riemannian metric by covering it with a locally finite family of charts
\((\psi_\alpha,O_\alpha)\) with \(\overline{O_\alpha}\) compact, defining a Riemannian
metric \((g_\alpha)_{ab}\) on each local coordinate neighborhood, and setting
\(g_{ab}=\sum_\alpha f_\alpha(g_\alpha)_{ab}\), where \(\{f_\alpha\}\) is a partition
of unity subordinate to \(\{O_\alpha\}\). Similarly, as discussed in Appendix~B, the
existence of a partition of unity allows us to define integration over a paracompact
manifold.
